%! Graphing Quadratic Functions & Transformations — Unit 3, Lesson 3.1 — Warm-Up (Answer Key)
\documentclass[10pt]{article}
\usepackage{algebra2-article}
\usepackage{algebra2-key}   % pulls in -boxes; do NOT also load -boxes
\newcommand{\TallMath}[1]{$\displaystyle #1\rule[-1.4em]{0pt}{3.2em}$}
\newcommand{\lgrid}[4]{%
  \draw[step=1, linegray!40, very thin] (#1,#3) grid (#2,#4);
  \draw[->, semithick, charcoal] (#1,0)--(#2,0) node[below right, font=\scriptsize]{$x$};
  \draw[->, semithick, charcoal] (0,#3)--(0,#4) node[left, font=\scriptsize]{$y$};}

\begin{document}

\pageheader{Unit 3, Lesson 3.1}{Warm-Up --- Answer Key}
\namedateperiod

\vspace{0.10in}

\begin{notesbox}{Getting Ready: Skills We'll Use Today}
\small These skills power today's lesson on graphing parabolas. Work quickly.

\vspace{4pt}
\textbf{1. The parent $y=x^2$.} Complete the table of the five ``anchor'' points, then read the graph.
\begin{center}
\begin{minipage}{0.50\linewidth}\centering
\renewcommand{\arraystretch}{1.3}
\begin{tabular}{|c|c|c|c|c|c|}
\hline
$x$ & $-2$ & $-1$ & $0$ & $1$ & $2$ \\ \hline
$y=x^2$ & \ans{$4$} & \ans{$1$} & \ans{$0$} & \ans{$1$} & \ans{$4$} \\ \hline
\end{tabular}
\end{minipage}%
\begin{minipage}{0.40\linewidth}\centering
\begin{tikzpicture}[scale=0.42]
  \lgrid{-3}{3}{-1}{5}
  \begin{scope}
    \clip (-3,-1) rectangle (3,5);
    \draw[very thick, forest, domain=-2.2:2.2, samples=60, smooth] plot (\x,{(\x)*(\x)});
  \end{scope}
  \fill[charcoal] (0,0) circle (2.5pt);
\end{tikzpicture}
\end{minipage}
\end{center}
The vertex of $y=x^2$ is the point \ans{$(0,0)$}, and it opens \ans{up}.

\vspace{4pt}
\textbf{2. Move one point (Unit 2 review).} The point $(2,4)$ is on the parent $y=x^2$. Give the new
location of \emph{that} point after each change:
\[
  y=x^2+3:\ \ans{$(2,7)$} \qquad y=(x-1)^2:\ \ans{$(3,4)$} \qquad y=-x^2:\ \ans{$(2,-4)$}
\]

\vspace{4pt}
\textbf{3. Expand a square.} Multiply out and simplify (write in the form $x^2+bx+c$).
\[
  (x-1)^2 - 4 = \ans{$x^2 - 2x - 3$}
\]
\end{notesbox}

\begin{teachernote}
Item~1 re-establishes the five anchor points every transformation moves. Item~2 previews the three
moves in point form: $+3$ raises $y$ by $3$; $(x-1)^2$ shifts the input \emph{right} 1 so the point at
$x=2$ now lives at $x=3$ (the sign-flip to flag today); $-x^2$ reflects over the $x$-axis. Item~3 is the
lesson's spine --- $(x-1)^2-4$ expands to $x^2-2x-3$, so vertex form and standard form describe the
\emph{same} parabola from Lesson~3.0. Keep the expansion on the board all period.
\end{teachernote}

\end{document}
